\section{Chapitre~\ref{sec:eetypes}. Les neurones comme composants}

Les modes de fonctionnement des cellules individuelles sont mesurés directement, par un courant injecté par électrode et un enregistrement de la tension; les mathématiques de l'intégrateur et de la division en classes relèvent de la théorie classique, tandis que l'usage par le cerveau de la reconfiguration des modes pour calculer reste une hypothèse.

{\footnotesize
\setlength{\tabcolsep}{3pt}\renewcommand{\arraystretch}{1.15}
\begin{longtable}{>{\raggedright}p{0.5cm}>{\raggedright}p{4.0cm}>{\raggedright}p{5.6cm}>{\raggedright}p{2.3cm}>{\raggedright\arraybackslash}p{1.7cm}}
\toprule
N\textsuperscript{o} & Affirmation & Expérience et ce qui est mesuré & Source & Statut \\
\midrule\endfirsthead
\toprule
N\textsuperscript{o} & Affirmation & Expérience et ce qui est mesuré & Source & Statut \\
\midrule\endhead
1 & Résonateur: un courant lent qui s'oppose aux variations de tension fait du neurone un filtre passe-bande dont le maximum se situe entre quelques hertz et quelques dizaines de hertz (section~\ref{sec:resonator}) & En tranches, on injectait dans des neurones \nt{GLUT} un courant sinusoïdal de fréquence croissante et on mesurait l'impédance $|Z(f)|$: maximum à $2$--$5$~Hz à $33^\circ$C (environ $7$~Hz à $38^\circ$C); la résonance est créée par des courants lents potassique et cationique, et amplifiée par le courant sodique persistant. Une revue décrit la même méthode pour de nombreuses classes de cellules & \cite{HuVervaekeStorm2002,HutcheonYarom2000} & mesuré \\
2 & Le courant lent se comporte comme une inductance; avec la capacité de la membrane, il forme un circuit oscillant & La réponse sous le seuil d'un axone à un courant a été mesurée et décrite par des équations de conductance linéarisées; le schéma équivalent contient une inductance \og phénoménologique\fg{} dont la valeur dépend de la tension & \cite{Mauro1970} & théorie \\
3 & Constante de temps d'un groupe à rétroaction $\tau_{\text{eff}}=\tau/(1-w)$~\eqref{eq:perfectint}; pour $\tau=0.1$~s et $w=0.995$, $20$~s & Déduite dans le chapitre de l'équation linéaire du groupe; proposée comme mécanisme de maintien de la position de l'œil dans un travail théorique & démonstration dans le chapitre; \cite{Seung1996} & théorie \\
4 & L'intégrateur de position de l'œil maintient son entrée grâce au réseau, et non grâce à une propriété d'une cellule isolée & Enregistrement intracellulaire chez le poisson, dans le noyau qui maintient la position de l'œil: après une rotation rapide, la fréquence du neurone change par palier et se maintient; un bref courant injecté dans une seule cellule ne provoque qu'un décalage passager, tandis que les paliers de tension persistent même sous le seuil: ce sont les entrées synaptiques qui les entretiennent & \cite{Aksay2001} & mesuré \\
5 & L'intégrateur sans fuite a besoin d'un réglage permanent par apprentissage; déréglé, il oublie ou part en saturation & On entraînait des poissons avec une rétroaction visuelle proportionnelle à $\pm$ la position de l'œil: le maintien devenait instable ou fuyant, et la constante de temps des neurones chutait de plus d'un ordre de grandeur, jusqu'à $<1$~s; la rétroaction visuelle normale rétablissait peu à peu la stabilité & \cite{Major2004} & mesuré \\
6 & Neurone bistable: une brève entrée enclenche un plateau durable, l'inhibition l'arrête; la propriété est activée par la sérotonine (section~\ref{sec:bistable}) & Enregistrement intracellulaire chez le chat de neurones \nt{ACET} qui commandent les muscles: une brève excitation synaptique fait passer le neurone en activité durable, une brève inhibition le remet à zéro; chez le chat spinal, l'état apparaît après administration d'un précurseur de la sérotonine & \cite{Hounsgaard1988} & mesuré \\
7 & Deux classes: les intégrateurs (fréquence croissant à partir de zéro) et les résonateurs (fréquence finie d'emblée au seuil) & Classes déduites de la théorie des bifurcations. Expérience: en tranches de cortex, les neurones \nt{GLUT} se mettent à décharger à une fréquence aussi basse qu'on veut (type~1), les neurones \nt{GABA} rapides démarrent d'un bond à fréquence finie (type~2) & \cite{Izhikevich2007,Tateno2004} & mesuré \\
8 & Un même neurone est intégrateur ou détecteur de coïncidences: l'inhibition raccourcit la fenêtre de sommation & En tranches d'hippocampe, la fenêtre dans laquelle les entrées excitatrices s'additionnent jusqu'au seuil est inférieure à $2$~ms; elle est raccourcie par l'inhibition directe qui suit immédiatement l'excitation; dans les dendrites, où l'inhibition est plus faible, la fenêtre est plus large & \cite{Pouille2001} & mesuré \\
9 & Ligne à retard faite d'axones (tableau~\ref{tab:eetypes}) & Chez la chouette effraie, on a mesuré les retards dans les axones qui vont aux détecteurs de coïncidences: des longueurs d'axone différentes donnent des retards différents, et les détecteurs sont accordés sur la différence des temps d'arrivée du son & \cite{Carr1990} & mesuré \\
10 & Stimulateur pendant la veille, cellule silencieuse pendant le sommeil & Les neurones \nt{SERO} déchargent de façon presque régulière le jour, ralentissent en sommeil lent et se taisent presque en sommeil paradoxal (détails au chapitre~\ref{sec:sleep}) & \cite{JacobsAzmitia1992,McGintyHarper1976} & mesuré \\
11 & Le composant est fixé par les canaux ioniques et par les canaux de diffusion qui les ajustent (proposition~\ref{prop:eetypes}) & Montré sur de petits réseaux: les substances régulatrices modifient les propriétés propres des neurones et la force des synapses, si bien qu'un même circuit produit des rythmes différents & \cite{Marder2012} & mesuré \\
12 & Le cerveau exploite la reconfiguration des modes pour calculer (proposition~\ref{prop:eetypes}) & Il n'existe aucune expérience directe où la reconfiguration du mode d'une cellule serait nécessaire pour résoudre une tâche; seulement des expériences où les régulateurs modifient la sortie du circuit & \cite{Marder2012} & hypothèse \\
\bottomrule
\end{longtable}}
